% Part: lambda-calculus
% Chapter: introduction
% Section: primitive-recursion

\documentclass[../../../include/open-logic-section]{subfiles}

\begin{document}

\olfileid{lam}{int}{pr}
\olsection{\usetoken{S}{lambda definable} વિધેયો આદિમ પુનરાવર્તન હેઠળ બંધ છે}

આદિમ પુનરાવર્તનની વાત આવે ત્યારે આખરે થોડું કામ કરવું પડે છે. આપણે
તબક્કાવાર આગળ વધશું. અગાઉની જેમ, ધારો કે આપણી પાસે એવા પદો~$G'$ અને
$H'$ છે જે અનુક્રમે વિધેયો~$g$ અને~$h$ને !!{lambda define} કરે છે.
આપણે નીચે પ્રમાણે વ્યાખ્યાયિત વિધેય~$f$ને !!{lambda define} કરતું
પદ~$F$ જોઈએ છે:
\begin{align*}
f(0, \vec z) & = g(\vec z) \\
f(x+1, \vec z) & = h(x, f(x,\vec z), \vec z).
\end{align*}
તેથી, સામાન્ય રીતે, આપેલાં લૅમ્બડા પદો~$G'$ અને~$H'$ માટે દરેક
પ્રાકૃતિક સંખ્યા~$n$ માટે નીચેની શરતો પૂરું કરતું પદ~$F$ શોધવું પૂરતું છે:
\begin{align*}
F(\num{0}, \vec z) & \equiv G'(\vec z) \\
F(\overline{n+1}, \vec z) & \equiv H'(\num{n}, F(\num{n}, \vec z), \vec z)
\end{align*}
કારણ કે~$G'$ અને~$H'$ અનુક્રમે~$g$ અને~$h$ને !!{lambda define} કરે
છે, $\vec z$ની જગ્યાએ અંકો~$\num{\vec m}$ મૂકતાં
$F(\num{n+1}, \num{\vec m})$ યોગ્ય ઉત્તરનાં સામાન્ય સ્વરૂપમાં
પહોંચશે.
\sourcecorrection{OLLAM-006}{સ્થિર અંગ્રેજી મૂળ શરૂઆતમાં~$G$
અને~$H$ આપે છે છતાં આગળ~$G'$ અને~$H'$ વાપરે છે, અને શોધવાનું પદ
ફરી~$H$ કહે છે. અહીં આપેલાં પદો~$G'$, $H'$ અને શોધવાનું પદ~$F$
રાખ્યાં છે; પહેલી બે પુનરાવર્તન-શરતોમાં પણ એ જ આપેલાં પદો છે.}
\sourcecorrection{OLLAM-007}{સ્થિર મૂળના પગથિયાના સમીકરણમાં
$h(z,f(x,\vec z),\vec z)$ લખાયું છે. અહીં પુનરાવર્તનનો સૂચક~$x$
છે, તેથી~$h(x,f(x,\vec z),\vec z)$ મૂક્યું છે.}

આ માટે, પરંતુ, નીચેની શરતો પૂરું કરતું પદ~$F$ શોધવું પૂરતું છે:
\begin{align*}
F(\num 0) & \equiv G \\
F(\overline {n+1}) & \equiv H(\num{n},F(\num{n}))
\intertext{દરેક પ્રાકૃતિક સંખ્યા $n$ માટે, જ્યાં}
G & = \lambd[\vec z][G'(\vec z)] \text{ અને}\\
H(u,v) & = \lambd[\vec z][H'(u,v(\vec z),\vec z)].
\end{align*}
બીજા શબ્દોમાં, લૅમ્બડા સંકેતનની યુક્તિથી વધારાના પ્રાચલો~$\vec z$
વિશે અલગથી ચિંતા કરવાની જરૂર રહેતી નથી---તે લૅમ્બડા પદમાં જ
સમાઈ જાય છે.
\sourcecorrection{OLLAM-008}{સ્થિર અંગ્રેજી મૂળમાં
$H(u,v)$ની વ્યાખ્યામાં~$v(u,\vec z)$ છે. અહીં~$v=F(\num{n})$
પહેલેથી~$n$ પર લાગુ કરેલું પદ છે અને હવે તે માત્ર~$\vec z$
સ્વીકારે છે; તેથી~$v(\vec z)$ મૂક્યું છે.}

પદ~$F$ વ્યાખ્યાયિત કરતાં પહેલાં ક્રમિત જોડીઓ સંભાળવાની એક રીત
જોઈએ. આગલું લેમ્મા તે આપે છે.

\begin{lem}
એવું લૅમ્બડા પદ~$D$ છે કે લૅમ્બડા પદોની દરેક જોડી~$M$ અને~$N$
માટે $D(M,N)(\num{0}) \red M$ અને $D(M,N)(\num{1}) \red N$.
\end{lem}

\begin{proof}
પહેલાં લૅમ્બડા પદ~$K$ને આ રીતે વ્યાખ્યાયિત કરો:
\[
K(y) = \lambd[x][y].
\]
બીજા શબ્દોમાં, $K$ એ પદ~$\lambd[y][\lambd[x][y]]$ છે. બીજી
દૃષ્ટિએ, દરેક~$M$ માટે $K(M)$ એવું અચલ વિધેય છે જે કોઈપણ ઇનપુટ
પર~$M$ પરત કરે છે.

હવે $D(x,y,z)$ને $D(x,y,z) = z (K(y))x$ વડે વ્યાખ્યાયિત કરો.
ત્યારે મળે છે:
\begin{align*}
D(M,N,\num 0) & \red \num 0 (K(N)) M \red M \text{ અને}\\
D(M,N,\num 1) & \red \num 1 (K(N)) M \red K(N) M \red N,
\end{align*}
જે જોઈએ હતું તે જ.
\end{proof}

વિચાર એ છે કે $D(M,N)$ જોડી~$\tuple{M, N}$નું નિરૂપણ કરે છે.
જો~$P$ આવી જોડીનું નિરૂપણ કરે એમ માનીએ, તો $P(\num 0)$ અને
$P(\num 1)$ અનુક્રમે ડાબા અને જમણા પ્રક્ષેપો, $(P)_0$ અને
$(P)_1$, દર્શાવે છે. આગળ આપણે આ છેલ્લાં સંકેતો વાપરીશું.

\begin{lem}
!!{lambda definable} વિધેયો આદિમ પુનરાવર્તન હેઠળ બંધ છે.
\end{lem}

\begin{proof}
આપેલાં કોઈપણ પદો~$G$ અને~$H$ માટે નીચેની શરતો પૂરું કરતું
પદ~$F$ શોધી શકાય એમ બતાવવાનું છે:
\begin{align*}
F(\num{0}) & \equiv G \\
F(\num{n+1}) & \equiv H(\num{n}, F(\num{n}))
\end{align*}
દરેક પ્રાકૃતિક સંખ્યા~$n$ માટે. મૂળ વિચાર અંકોને પુનરાવર્તક તરીકે
વાપરી \emph{જોડીઓ}ની નીચેની શ્રેણી ગણવાનો છે:
\[
\tuple{\num 0, F(\num 0)}, \tuple{\num 1, F(\num 1)}, \dots,
\]
પ્રથમ જોડી માત્ર~$\tuple{\num 0, G}$ છે તે ધ્યાનમાં લો. જોડી
$\tuple{\num{n}, F(\num{n})}$ આપેલી હોય, તો આગલી જોડી
$\tuple{\num{n+1}, F(\num{n+1})}$ એ
$\tuple{\num{n+1}, H(\num{n}, F(\num{n}))}$ને સમકક્ષ હોવી
જોઈએ. આ એક પગથિયાનો ફેરફાર કરતું લૅમ્બડા પદ~$T$ રચીશું.

વિગતો આ પ્રમાણે છે. $T(u)$ને આ રીતે વ્યાખ્યાયિત કરો:
\[
T(u) = \tuple{S((u)_0), H((u)_0,(u)_1)}.
\]
અહીં પ્રથમ ઘટકમાં અગાઉ વ્યાખ્યાયિત ઉત્તરગામી લૅમ્બડા પદ વપરાય છે.
હવે ચકાસવું સરળ છે કે કોઈપણ સંખ્યા~$n$ માટે
\[
T(\tuple{\num{n}, M}) \red \tuple{\num{n+1}, H(\num{n}, M)}.
\]
ઉપર સૂચવ્યા મુજબ, $G$ અને~$H$ આપેલાં હોય ત્યારે~$F(u)$ને
આ રીતે વ્યાખ્યાયિત કરો:
\[
F(u) = (u(T,\tuple{\num 0, G}))_1.
\]
બીજા શબ્દોમાં, ઇનપુટ~$\num{n}$ પર~$F$, જોડી
$\tuple{\num 0, G}$ પરથી~$T$ને $n$ વખત લાગુ કરે છે અને પછી
બીજો ઘટક પરત કરે છે. શરૂઆતમાં મળે છે:
\begin{enumerate}
\item $\num{0} (T, \tuple{\num 0, G}) \equiv \tuple{\num{0}, G}$
\item $F(\num{0}) \equiv G$
\end{enumerate}
$n$ પર આગમનથી દરેક પ્રાકૃતિક સંખ્યા માટે નીચેની બે બાબતો
બતાવી શકીએ:
\begin{enumerate}
\item $\num{n+1}(T, \tuple{\num{0}, G}) \equiv \tuple{\num{n+1}, F(\num{n+1})}$
\item $F(\num{n+1}) \equiv H(\num{n}, F(\num{n}))$
\end{enumerate}
બીજી બાબત માટે,
\begin{align*}
F(\num{n+1}) & \red (\num{n+1}(T, \tuple{\num 0, G}))_1 \\
& \equiv (T(\num{n} (T, \tuple{\num 0, G})))_1 \\
& \equiv (T(\tuple{\num{n}, F(\num{n})}))_1 \\
& \equiv (\tuple{\num{n+1}, H(\num{n}, F(\num{n}))})_1 \\
& \equiv H(\num{n}, F(\num{n})).
\end{align*}
અહીં છેલ્લેથી બીજી પંક્તિમાં આગમનની ધારણા વાપરી છે. પહેલી
બાબત માટે,
\begin{align*}
\num{n+1} (T, \tuple{\num 0, G}) &
\equiv T(\num{n} (T, \tuple{\num 0, G})) \\
& \equiv T( \tuple{\num{n}, F(\num{n})}) \\
& \equiv \tuple{\num{n+1}, H(\num{n}, F(\num{n}))} \\
& \equiv \tuple{\num{n+1}, F(\num{n+1})}.
\end{align*}
અહીં છેલ્લી પંક્તિમાં બીજી બાબત વાપરી છે. આમ આપણે બતાવ્યું કે
$F(\num 0) \equiv G$ અને દરેક~$n$ માટે
$F(\num {n+1}) \equiv H(\num n, F(\num{n}))$; બરાબર એ જ
આપણને જોઈએ હતું.
\end{proof}

\end{document}
